% Numeric source: paper/shared/evidence/raw/iclr_model_behavior/analysis_v1/results.json
\begin{abstract}
The same scene need not produce the same decisions. Converting its materials
preserves geometry but changes the visual evidence available to an embodied
model. In \NavN\ paired navigation episodes, success falls from
\NavSrOriginal\% to \NavSrConverted\%, yet \NavFlips\ episodes change success
status: gains and losses largely cancel in the average. We investigate this
hidden turnover through paired visual grounding and navigation experiments.
A readout study across six vision-language model configurations separates
localization from response-format and coordinate-scale effects. Controlled
two-view tests then require models to follow a queried object across positions,
and target-window crossovers reveal that changes to the target and its
surroundings can shift grounding in opposite directions. In navigation,
17 previously successful episodes fail while 13 previously failed episodes
succeed under converted materials. These results put the identity of successful
decisions alongside mean accuracy as a central object of evaluation. Material
representation is therefore more than an asset-format choice: it is an input
variable through which we can study the visual dependence of embodied models.
\end{abstract}
